\section{Innovation II: Technology}

\subsection*{Technology \& Industry Evolution}
\begin{itemize}
    \item \textbf{Kondratieff Waves:} Successive waves of new technologies (e.g., Industrial
          Revolution, Steam, Electricity, Mass Production, Information Age) drive periods
          of economic growth.
    \item \textbf{Abernathy \& Utterback (1978) Model:} Industry evolution follows three distinct
          phases:
          \begin{enumerate}
              \item \textbf{Fluid Phase (Stage 1):} Characterized by high product innovation, exploration,
                    uncertainty, and flexibility.
              \item \textbf{Transitional Phase (Stage 2):} Marked by the emergence of a \textit{dominant design}—an accepted product architecture that defines the look and
                    functionality for the industry.
              \item \textbf{Specific Phase (Stage 3):} Focus shifts toward high process innovation,
                    emphasizing standardization and integration.
          \end{enumerate}
\end{itemize}

\subsection*{Market Leadership \& Competence}
\begin{itemize}
    \item \textbf{Competence Enhancing Innovation:} Builds directly on a firm's existing
          knowledge base (e.g., upgrading from Pentium III to Pentium 4). This typically
          results in no change in market leadership.
    \item \textbf{Competence Destroying Innovation:} Renders a firm's existing knowledge base
          obsolete (e.g., mobile phone experts competing against the iPhone's computer
          science architecture). This often leads to a shift in market leadership.
\end{itemize}

\subsection*{Demand-Side Innovation \& Disruption}
\begin{itemize}
    \item \textbf{Demand vs. Technology:} Research shows 75\% of successful innovations stem from
          understanding customer needs, while only 25\% rely purely on technological
          opportunity.
    \item \textbf{Disruptive Innovation (e.g., Tinder vs. Match.com):}
          \begin{enumerate}
              \item \textbf{Sustaining Innovations:} Incumbents focus investments on dimensions that
                    mainstream customers value.
              \item \textbf{Niche Entry:} New entrants focus on overlooked dimensions (e.g., gamified
                    mobile swiping) appreciated by a niche or simpler segment.
              \item \textbf{Displacement:} As the niche expands, the new entrant can displace the
                    established incumbent.
          \end{enumerate}
\end{itemize}